\section{Method}\label{sec:method}

\subsection{The mapping (prior work, condition C)}\label{sec:heo}
We reproduce the recipe of~\citep{heo2026crossmodel} as written.
\begin{itemize}
\item \emph{Features.} Per source layer, all KV heads are flattened. Keys have the source RoPE removed; values are used as is.
\item \emph{Targets.} The destination's own keys (RoPE removed) and values, per layer and head.
\item \emph{Regression.} Centred ridge with $\lambda=0.01$ and bias $b=\bar y-\bar xW$. Covariance is accumulated in float32 and solved in float64.
\item \emph{Data.} 500 FineWeb-Edu~\citep{penedo2024fineweb} sequences of 1\,024 tokens, sub-sampled with stride 4.
\item \emph{Source selection.} For each destination layer, the top-$k$ source layers are kept by a fixed ranking: head-averaged single-source $R^2$, averaged over K and V.
\end{itemize}
Keys are re-rotated with the destination RoPE at the same absolute positions.

We had to make two interpretations, and we state both.
(i) The text writes $X^\top X+\lambda I$ without saying whether $X^\top X$ is normalised.
On $\sim$128\,K rows an unnormalised Gram matrix would make $\lambda=0.01$ a no-op, and their reported collapse at $\lambda=1$ could not happen.
We therefore read $X^\top X$ as the centred covariance and also sweep $\lambda$ (\cref{tab:lambda}).
(ii) Heo et al.\ choose $k$ by downstream accuracy.
We choose $k$ by held-out prefix-KV negative log-likelihood (NLL), measured on documents disjoint from both the calibration and the evaluation splits.

Implementation checks in \texttt{fit.selftest} and \texttt{migrate.selftest}:
a planted linear map is recovered to $R^2=1$;
the $R^2$ computed from sufficient statistics equals the $R^2$ of actually applying the fitted mapper to within $10^{-7}$;
and with an identity mapper, transplanting reproduces a native re-prefill to within $1.4\times10^{-6}$.

\subsection{Mid-decode transplant}\label{sec:transplant}
The decode loop keeps one invariant: the cache holds KV for exactly the tokens fed so far, and \texttt{pending} is the token already shown to the user but not yet fed.
At a switch the receiver gets the mapped cache and feeds \texttt{pending} itself, so no shown token is retracted.
Molt adds two optional \emph{anchors}, which we treat as engineering choices and not as contributions.
\begin{itemize}
\item The first $s$ positions are replaced by the receiver's native KV. This is exact, because causal positions depend only on earlier tokens; it follows the sink observation of~\citep{xiao2023streamingllm}.
\item The last $t$ positions are re-fed natively over the mapped history.
\end{itemize}
Setting $t=\text{all}$ gives re-prefill (condition A$'$).
$s$ and $t\le 32$ are chosen on held-out text only.

\subsection{The memory order (the claim)}\label{sec:order}
Releasing a model on MLX is not one call.
We measured three facts (\cref{app:mlx}).
(1) After \texttt{del}, buffers stay in MLX's allocator cache and are not reused for arrays of other sizes.
(2) \texttt{mx.clear\_cache()} right after \texttt{del} can free nothing, because a buffer returns to the cache only after the GPU has finished the last command buffer that used it.
(3) After synchronising and clearing, the OS still reclaims the memory asynchronously, in about $10^2$\,ms.
The \emph{ordered} schedule therefore runs as follows:
\begin{enumerate}
\item snapshot the source KV;
\item drop the source weights, synchronise, clear, and poll phys\_footprint until it has fallen by 90\% of the source weight bytes;
\item load the mapper and map;
\item drop the mapper and source KV the same way;
\item load the receiver;
\item apply the anchors and resume.
\end{enumerate}
\Cref{tab:schedules} lists the three alternatives we measure against it.
\emph{colocated} is the prefill hand-off shape.
\emph{del-only} is the natural ``drop it, then load''.
\emph{clear-no-wait} is the call a developer would write.

\paragraph{The trigger.}
Before every token the loop reads the headroom.
On iOS that is the return value of \texttt{os\_proc\_available\_memory()}; on the Mac it is $\text{cap}-\text{phys\_footprint}$.
It migrates when the headroom falls below what \emph{the chosen schedule} needs plus the next demand chunk.
Each schedule thus gets its own best trigger, and a co-resident schedule is never killed merely for reacting late.
